\documentclass{article}
\usepackage{kotex}
\usepackage{setspace}  
\setstretch{1.15}
\usepackage[a4paper, 
            left = 2.5cm,
            right = 2.5cm,
            top = 2cm,
            bottom = 2cm]{geometry}
\usepackage{fancyhdr}                 
\pagestyle{fancy}
\fancyhf{}                             
\fancyfoot[L]{\nouppercase{통사론 (2025학년도 2학기)}} 
\fancyfoot[R]{\thepage}                
\renewcommand{\headrulewidth}{0pt}    
\renewcommand{\footrulewidth}{0.4pt}   

\begin{document}
\large

\title{통사론 중간고사 (2025학년도 2학기)}
\author{교수자: 강초롱}
\date{2025 Fall}
\maketitle
\thispagestyle{fancy}
\bigskip

\section{Consider the following data taken from Turkish \normalsize(3 points)}
\medskip \noindent 1) Siz \textbf{gid}-iyor.sun.uz.\par 
\noindent 2) Ali \"{u}\c{c} ye\c{s}il top \textbf{al}-d\i,\par
\noindent 3) Senin s\i{}nav\i{} \textbf{ge\c{c}}-meni \textbf{isti}-yorum.\par
\bigskip
\noindent a. In the data, PRES means ``present tense'', PST means ``past tense''. 
What part of speech are `gid', `al', `ge\c{c}', and `isti'? 
What tests did you use to figure out what speech they were? \par
\vspace{1cm}
\noindent b. Provide the PSRs needed to generate the trees (assume that every clause is a CP, and \textit{\"{u}\c{c}} ``three'' is an adjective). \par
\vspace{1cm}

\section{Consider the following sentence \normalsize}
\textit{The teacher saw that her students fought constantly after school in the parking lot.} \par
\bigskip
\noindent a. For the following words, indicate their part of speech AND \par
\medskip \noindent i) if the word is an open class part of speech, explain the (morphological/syntactic)\\ distributional criteria by which you came up with that classification.\par
\noindent ii) if the word is closed part of speech, indicate it with ``closed''.\par
\noindent iii) Mark whether each word is a lexical or functional part of speech. (4 points) \par
\begin{description}
    \item[\textbf{The}]
    \item[\textbf{Teacher}] 
    \item[\textbf{Saw}] 
    \item[\textbf{That}] 
    \item[\textbf{Her}] 
    \item[\textbf{Constantly}]
    \item[\textbf{After}]
    \item[\textbf{Parking}]   
\end{description}

\newpage
\noindent b. The above sentence is ambiguous (assume that the ambiguity comes from the phrase ``in the parkng lot''). Draw the two PSR trees that distinguish the two possible meanings of the sentence. Provide a Korean paraphrase for each meaning. (2 points)

\newpage 
\noindent c. Draw the two X-bar theoretic trees that distinguish the two possible meanings of the sentence. Provide a Korean paraphrase for each meaning. (2 points)

\newpage
\noindent d. Answer the following questions (6 points) \par
\medskip \noindent i) Can the two theories (PSRs and X-bar theory) explain the ambiguity? If so, explain why. You must refer to the tree structures you provided above in your answer. \par
\vspace{1.5cm}
\noindent ii) In the slightly modified sentence below, the interpretation in which \textit{the parents fought after school at the park and the teacher saw it} is available. Can both theories account for this reading? \par
\medskip \noindent \textit{The teacher saw that her students fought constantly after school in the parking lot. She saw that their parents did so at the park as well.}\par
\vspace{1.5cm}
\noindent iii) Explain why (a) is ungrammatical but (b) is grammatical, referring to the rules of X-bar theory.\par
\medskip \noindent (a) \*The teacher saw \textbf{the fight} that her students fought constantly after school in the parking lot.\par
\medskip \noindent (b) The teacher saw that her students fought constantly after school in the parking lot \textbf{every day}.

\newpage
\section{Consider the following data from Breton, a Celtic language. \normalsize{(7 points)}}

\newpage
\section{Draw the X-bar theoretic tree for the following sentence, and determine whether the matrix clause and the embedded clause are finite or non-finite. You must provide evidence for your answer. \normalsize{(3 points)}}
\medskip
\begin{center} 
    \textbf{I asked him to leave.}
\end{center}

\end{document}